\begin{claim}{D.2.56}{Proved}{The next two homology groups of the connected cover}
Let $L=Y\langle7\rangle$, $E=\Omega^2L$, and retain the
double-loop fiber $F$ of D.2.18.  All groups in this claim have
coefficients $\FF_2$.  The edge map for
$E\to F\xrightarrow{p}K(\FF_2,4)$ identifies
\[
 H^{\le12}E=
 \left(H^*F/(p^*\widetilde H^*K(\FF_2,4))H^*F\right)^{\le12}.
 \tag{D.2.56a}
\]
In particular $\dim H_{11}E=20$ and $\dim H_{12}E=32$.
Use the previously specified classes
\[
 \begin{gathered}
 u=\xi a\in H_5E,\quad D=\{a,b\},\quad E_6=b^2+p_3\in H_6E,\\
 S=\xi b+P+Q,\quad R=\{a,c\}\in H_7E.
 \end{gathered}
 \tag{D.2.56b}
\]
Actual double homology suspension induces the quotient isomorphisms
\[
 \begin{aligned}
 H_{13}L&\cong H_{11}E/\langle uD,uE_6,\xi u\rangle,
                 &\dim H_{13}L&=17,\\
 H_{14}L&\cong H_{12}E/
 \langle uS,uR,D^2,DE_6,E_6^2,\{u,D\},\{u,E_6\}\rangle,
                 &\dim H_{14}L&=25.
 \end{aligned}
 \tag{D.2.56c}
\]
These descriptions commute with Steenrod operations whenever both
degrees are in the displayed range.  They identify the raw homology
inputs for the potential incoming sixth and seventh differentials
in D.2.54, but do not compute their later-page subquotients, higher
Bocksteins, or stable attaching maps.
\end{claim}
\begin{proof}
\textbf{A bounded Hopf-algebra argument.}
Put $A=H^*K(\FF_2,4)$, $C=H^*F$, and $A'=\operatorname{im}(p^*)$.
The map $p$ is the double loop of the degree-six Postnikov map of
the five-connected cover of $Y$.  It is therefore an H-map, so
$A'$ is a connected Hopf subalgebra of $C$.  Since cup multiplication
is commutative, Milnor--Moore's freeness theorem applies:
$C$ is a graded free $A'$-module
\cite[Proposition 4.9 and the following paragraph, pp.~229--230]
{I_MilnorMoore1965}.  The field hypothesis supplies all the required
graded vector-space splittings.

We will prove that $p^*$ is injective through degree thirteen.
Its kernel then has least possible degree fourteen.  A minimal
graded $A$-free resolution of $A'$ has first positive generators
in internal degree at least fourteen.  Subsequent minimal
differentials have entries of degree at least four, the least
positive degree of $A$.  Consequently
\[
 \operatorname{Tor}_s^A(A',\FF_2)_t=0
       \quad(s>0,\ t<14+4(s-1)).
 \tag{D.2.56d}
\]
The same bound holds for $\operatorname{Tor}_s^A(C,\FF_2)$,
because $C$ is a direct sum of nonnegative shifts of $A'$.
Positive Tor has total degree at least $13+3(s-1)$.
The convergent Eilenberg--Moore sequence for this simply connected
base thus has only filtration zero through total degree twelve,
and no incoming differential can affect that range.  Its edge
map is restriction, proving \textup{(D.2.56a)} once the promised
injectivity is established.  This does not assume global
injectivity of $p^*$.

\textbf{The three new primitive tests.}
The proof of the cover calculation in D.2.22 establishes
injectivity through degree eleven.  A nonzero element of least
degree in the kernel of a connected Hopf map is primitive.
The Serre polynomial generators of $A$ are primitive, and the
other primitives are their iterated squares.  The only new
primitive directions to test are
\[
 \begin{array}{c|l}
 12&(\operatorname{Sq}^2x_4)^2,
             \quad\operatorname{Sq}^5\operatorname{Sq}^2
                       \operatorname{Sq}^1x_4,\\
 13&\operatorname{Sq}^6\operatorname{Sq}^3x_4,
             \quad\operatorname{Sq}^6\operatorname{Sq}^2
                       \operatorname{Sq}^1x_4.
 \end{array}
 \tag{D.2.56e}
\]
Here $x_4$ is the fundamental class of $K(\FF_2,4)$.

Write $u_B,v_B,w_B,t_B$ for the degree $2,3,5,9$ Serre
generators of $H^*B$, and $[m]$ for their dual homology labels
in $T=\Omega^2\Sigma^2B$.  The formal label series
$\mathcal G=\sum_m[m]m$ is group-like.  Hence coefficients of
$\mathcal G^{-1}t_B\partial\mathcal G/\partial t_B$,
and the analogous $w_B$ logarithmic derivative, are primitive.
Let
\[
 \begin{aligned}
 Z_{11}&=[u_Bt_B]+[u_B][t_B],\\
 Z_{12}&=[v_Bt_B]+[v_B][t_B],\\
 Z_{13}^{(1)}&=[u_B^4w_B]
       (\mathcal G^{-1}w_B\partial\mathcal G/\partial w_B),\\
 Z_{13}^{(2)}&=[u_B^2t_B]
       (\mathcal G^{-1}t_B\partial\mathcal G/\partial t_B).
 \end{aligned}
 \tag{D.2.56f}
\]
In the last two lines the initial bracket means coefficient
extraction in the formal variables, not multiplication by a
homology label.  All four classes lie in $H_*F$: they are
primitive in $H_*T$, their images in $H_*B$ are zero since
$B$ has no primitive homology in degrees $11,12,13$, and
the actual coalgebra splitting $H_*T=H_*B\otimes H_*F$
splits primitives into those of the two factors.

Cartan and the Serre operations give
\[
 \operatorname{Sq}_*^1Z_{12}=Z_{11},\qquad
 \operatorname{Sq}_*^6Z_{13}^{(1)}=P,\qquad
 \operatorname{Sq}_*^6Z_{13}^{(2)}=Q.
 \tag{D.2.56g}
\]
For the last two identities the linear-label coefficients are
exactly those in
\[
 \operatorname{Sq}^6(u_B^2v_B)=u_B^4w_B,
 \qquad
 \operatorname{Sq}^6(u_Bw_B)=v_Bw_B^2+u_B^2t_B.
 \tag{D.2.56h}
\]
Primitivity then fixes the remaining terms, since the polynomial
primitive degree-seven subspace has basis $P,Q$.  The existing
degree-eleven test pairs $Z_{11}$ to one with
$p^*\operatorname{Sq}^4\operatorname{Sq}^2\operatorname{Sq}^1x_4$.
Thus $Z_{12}$ detects the new degree-twelve generator and pairs
to zero with the square in \textup{(D.2.56e)}.

The remaining square is detected by $[v_B^2]^2$.
Its reduced coproduct is $[v_B]^2\otimes[v_B]^2$.
Both factors, and $[v_B^2]^2$ itself, map to zero under
$T\to B$ by divided-power multiplication in $H_*B$.
Its full right $H_*B$-coaction is therefore
$[v_B^2]^2\otimes1$.  In the actual coalgebra splitting
$H_*B\otimes H_*F$, these coinvariants are precisely
$1\otimes H_*F$; applying the $B$-counit to the first
factor of the coinvariant identity proves this directly.
This justifies its fiber inclusion.  Since
$\operatorname{Sq}_*^2([v_B]^2)=a^2$, it pairs to one
with $(p^*\operatorname{Sq}^2x_4)^2$.

Finally $p^*\operatorname{Sq}^3x_4$ and
$p^*\operatorname{Sq}^2\operatorname{Sq}^1x_4$ pair on
$(P,Q)$ as $(1,0)$ and $(1,1)$.  Equations
\textup{(D.2.56g)--(D.2.56h)} give an invertible matrix
for the two degree-thirteen directions.  This proves
injectivity through thirteen and completes the edge argument.

\textbf{Dimensions and low products.}
The free restricted-Lie description of the double-loop homology
\cite[Chapter III, Theorems 3.1--3.2 and Lemma 3.8]
{I_CohenLadaMay1976}, the space splitting $T\simeq B\times F$,
and Hopf freeness give the following finite Hilbert-series counts:
\[
 \begin{array}{c|rrrrrrrr}
 i&5&6&7&8&9&10&11&12\\\hline
 \dim H_iE&1&2&2&5&8&12&20&32.
 \end{array}
 \tag{D.2.56i}
\]
For reproducibility, the ordinary free-Lie counts in shifted
degrees $3$ through $14$ are
$1,1,1,2,3,3,6,9,12,17,28,37$; adjoining their restricted
squares and taking the polynomial algebra gives
$\dim H_{11}T=69$, $\dim H_{12}T=112$.
Dividing first by the series of $H_*B$ gives $30,49$
for $F$, and then by that of $H^*K(\FF_2,4)$ gives
$20,32$ for $E$.  The exact finite certificate
\texttt{code/I\_cover\_extension.py} checks these integer
series, all four primitive coproducts and
\textup{(D.2.56g)}.  No homotopy calculation is inferred
from this certificate.

In the polynomial free-E2 homology of $T$, the two products
$uD,uE_6$, the five products
$uS,uR,D^2,DE_6,E_6^2$, and the additional generator
$\xi u$ have the claimed independence in their respective
degrees.  Indeed $u,D,R$ are distinct free generators,
$S$ has the distinct $\xi b$ term, and $E_6$ has the
distinct linear-label $f$ term.  Their triangular expressions
generate a polynomial subalgebra.  The class $\xi u$ is
a separate free generator.

\textbf{The two bar filtrations.}
Use the actual reduced bar constructions $BE\simeq\Omega L$
and $B(\Omega L)\simeq L$.  Their field-coefficient
filtration-one edges are homology suspension, and their
$E^1$ differentials multiply adjacent factors
\cite[Chapter 1, Theorem 1.5 and its footnote, p.~7]
{I_MadsenMilgram1979}.  In $BE$, factors have degree at
least five.  Length three starts in total degree eighteen.
In total degree twelve the only length-two term is $u|u$;
its differential is the nonzero class $u^2$.
The length-two boundaries into the degree-twelve
filtration-one term are $uD,uE_6$.  Therefore
\[
 H_{12}(\Omega L)=H_{11}E/\langle uD,uE_6\rangle,
                  \qquad\dim=18.
 \tag{D.2.56j}
\]
In the second bar construction the term
$\sigma u|\sigma u$ has differential $(\sigma u)^2$.
The top-operation suspension formula identifies this
nonzero class with $\sigma(\xi u)$
\cite[Chapter III, Theorem 1.4(2), with Theorem 1.1(2)]
{I_CohenLadaMay1976}.  This proves the first quotient in
\textup{(D.2.56c)}.

In total degree thirteen of $BE$, filtration one has
dimension $32-5=27$, after quotienting by the five
displayed degree-twelve products.  There are also the
two surviving length-two cycles
\[
 u|D+D|u,\qquad u|E_6+E_6|u.
 \tag{D.2.56k}
\]
They are the shuffle-leading terms of
$(\sigma u)(\sigma D)$ and $(\sigma u)(\sigma E_6)$.
The two reversed products must be retained: $\Omega L$
is only an E1 space.  Their differences have filtration
one and are exactly
\[
 \sigma\{u,D\},\qquad\sigma\{u,E_6\},
 \tag{D.2.56l}
\]
by the Browder suspension and commutator formulas
\cite[Chapter III, Theorem 1.4(3) and Theorem 1.2(2)]
{I_CohenLadaMay1976}.  They are independent modulo the
five products: their free-Lie indecomposable parts have
the respective nonzero words $[a,[a,[a,b]]]$ and
$[a,[a,f]]$, with different linear labels and arities.
Thus the four ordered products have rank four, with
two independent shuffle-leading terms and two independent
lower commutators.

The second bar differential kills all four products.
There are no other terms or higher incoming differentials
in this degree, since its factors have degree at least six
and length three starts in total degree twenty-one.
It follows that double suspension is onto in degree
fourteen and its kernel is exactly the seven-dimensional
subspace in \textup{(D.2.56c)}.  Its target has dimension
$27+2-4=25$.  Naturality of actual homology suspension
proves the stated operation compatibility.  The computation
ends with these ordinary homology quotients; no value of
a later stable differential has been used.
\end{proof}
